\pdfbookmark[0]{Anzeigen. Monge: Géométrie descriptive. 1813.}{anzeigen-monge-1813}
\begin{center}
\rule{.85\textwidth}{.3pt}\par
Göttingische gelehrte Anzeigen.\quad 31. Juli 1813.\par
\rule{.85\textwidth}{.3pt}
\end{center}
\hspace*{\parindent}\textit{Géométrie descriptive par} \textsc{Gaspard Monge}, \textit{de l’institut des sciences etc. Nouvelle édition. Avec un supplément par} \textsc{M. Hachette}, \textit{instituteur à l’école impériale polytechnique etc.} Paris, bei \textsc{J. Klostermann} dem Jüngeren. 162 und 118 Seiten im Quartformat.

Die Geometrie untersucht räumliche Verhältnisse. Sie gliedert sich in zwei große Gebiete, je nachdem, ob der Raum nur in zwei Dimensionen, also in der Ebene, oder in allen drei Dimensionen zugleich betrachtet wird. Es ist leicht einzusehen, dass das zweite Gebiet seiner Natur nach viel umfangreicher sein und weit mehr unterschiedliche Fragen und Untersuchungen bieten muss als das erste. Wenn also schon in unserer elementaren Geometrie die Geometrie der Ebene mehr Raum einnimmt als die Geometrie des Raumes, liegt das nur daran, dass diese verhältnismäßig viel weniger entwickelt und ausgearbeitet ist. Gerade Untersuchungen der zweiten Art hat man in neuerer Zeit tatsächlich lieber mit den Mitteln der Analysis behandelt und sie damit gewissermaßen der Geometrie entzogen, die sich allein auf die unmittelbare Anschauung stützt. Die Vorzüge der analytischen gegenüber der geometrischen Behandlung sind unbestreitbar. Ihre Kürze und Einfachheit, ihr gleichmäßiges Vorgehen und vor allem ihre Allgemeinheit treten gewöhnlich umso deutlicher hervor, je schwieriger und verwickelter die Untersuchungen sind. Gleichwohl bleibt es sehr wichtig, auch die geometrische Methode weiter zu pflegen. In manchen einzelnen %
